\documentclass[11pt,a4paper]{article}
\usepackage[utf8]{inputenc}
\usepackage[T1]{fontenc}
\usepackage{amsmath,amssymb,amsfonts}
\usepackage[margin=1in]{geometry}
\usepackage{enumitem}
\usepackage{hyperref}
\usepackage{fancyhdr}
\usepackage{titlesec}
\usepackage{tocloft}
\newtheorem{example}{Example}
\newtheorem{theorem}{Theorem}
\newtheorem{definition}{Definition}
\newtheorem{lemma}{Lemma}
\newtheorem{corollary}{Corollary}
\newtheorem{remark}{Remark}
\pagestyle{fancy}
\fancyhf{}
\fancyhead[L]{\small Lecture Notes: Techniques of Integration}
\fancyhead[R]{\small \thepage}
\renewcommand{\headrulewidth}{0.4pt}
\titleformat{\section}{\Large\bfseries}{Section \thesection}{1em}{}
\titleformat{\subsection}{\large\bfseries}{\thesubsection}{1em}{}
\renewcommand{\cftsecfont}{\bfseries}
\renewcommand{\cftsubsecfont}{\normalfont}
\begin{document}
\begin{titlepage}
\centering
\vspace*{4cm}
{\Huge\bfseries Lecture Notes:\\ Techniques of Integration\par}
\vspace{1cm}
{\Large Calculus I\par}
\vspace{2cm}
{\large \textbf{Author:} Bakdaulet Sotsial\par}
\vspace{1cm}
\vfill
{\small Astana IT University \par}
\vspace*{2cm}
\end{titlepage}
\tableofcontents
\newpage
\section{Integration by Parts}
\begin{theorem}[Integration by Parts Formula]
If $u = f(x)$ and $v = g(x)$ are differentiable functions, then
\[
\int u \, dv = uv - \int v \, du.
\]
In terms of definite integrals,
\[
\int_a^b u \, dv = \left[ uv \right]_a^b - \int_a^b v \, du.
\]
\end{theorem}
\begin{proof}
The product rule for differentiation states that
\[
\frac{d}{dx}(uv) = u \frac{dv}{dx} + v \frac{du}{dx}.
\]
Integrating both sides with respect to $x$ gives
\[
uv = \int u \, dv + \int v \, du.
\]
Rearranging yields $\int u \, dv = uv - \int v \, du$.
\end{proof}
\begin{remark}[Choosing $u$ and $dv$]
A useful guideline for choosing $u$ is the acronym LIATE: Logarithmic, Inverse trigonometric, Algebraic, Trigonometric, Exponential. The function that appears earlier in this list should generally be chosen as $u$, and the remaining part as $dv$.
\end{remark}
\begin{example}
Evaluate $\displaystyle \int x e^x \, dx$.
Let $u = x$ and $dv = e^x \, dx$. Then $du = dx$ and $v = e^x$. Applying the formula,
\[
\int x e^x \, dx = x e^x - \int e^x \, dx = x e^x - e^x + C = e^x(x - 1) + C.
\]
\end{example}
\begin{example}
Evaluate $\displaystyle \int \ln x \, dx$.
Let $u = \ln x$ and $dv = dx$. Then $du = \frac{1}{x} \, dx$ and $v = x$. Thus,
\[
\int \ln x \, dx = x \ln x - \int x \cdot \frac{1}{x} \, dx = x \ln x - \int 1 \, dx = x \ln x - x + C.
\]
\end{example}
\begin{example}
Evaluate $\displaystyle \int x \sin x \, dx$.
Let $u = x$ and $dv = \sin x \, dx$. Then $du = dx$ and $v = -\cos x$. Thus,
\[
\int x \sin x \, dx = -x \cos x - \int (-\cos x) \, dx = -x \cos x + \int \cos x \, dx = -x \cos x + \sin x + C.
\]
\end{example}
\begin{example}
Evaluate $\displaystyle \int x^2 e^x \, dx$.
Apply integration by parts twice. First, let $u = x^2$ and $dv = e^x \, dx$. Then $du = 2x \, dx$ and $v = e^x$:
\[
\int x^2 e^x \, dx = x^2 e^x - 2 \int x e^x \, dx.
\]
From the previous example, $\int x e^x \, dx = e^x(x - 1) + C$. Therefore,
\[
\int x^2 e^x \, dx = x^2 e^x - 2e^x(x - 1) + C = e^x(x^2 - 2x + 2) + C.
\]
\end{example}
\begin{example}
Evaluate $\displaystyle \int e^x \sin x \, dx$.
Let $I = \int e^x \sin x \, dx$. Apply integration by parts with $u = \sin x$ and $dv = e^x \, dx$. Then $du = \cos x \, dx$ and $v = e^x$:
\[
I = e^x \sin x - \int e^x \cos x \, dx.
\]
Apply integration by parts again to $\int e^x \cos x \, dx$ with $u = \cos x$ and $dv = e^x \, dx$. Then $du = -\sin x \, dx$ and $v = e^x$:
\[
\int e^x \cos x \, dx = e^x \cos x + \int e^x \sin x \, dx = e^x \cos x + I.
\]
Substituting back,
\[
I = e^x \sin x - (e^x \cos x + I) = e^x \sin x - e^x \cos x - I.
\]
Thus, $2I = e^x(\sin x - \cos x)$, so
\[
I = \frac{e^x}{2}(\sin x - \cos x) + C.
\]
\end{example}
\begin{example}
Evaluate $\displaystyle \int \arctan x \, dx$.
Let $u = \arctan x$ and $dv = dx$. Then $du = \frac{1}{1+x^2} \, dx$ and $v = x$. Thus,
\[
\int \arctan x \, dx = x \arctan x - \int \frac{x}{1+x^2} \, dx.
\]
For the remaining integral, let $w = 1 + x^2$, so $dw = 2x \, dx$:
\[
\int \frac{x}{1+x^2} \, dx = \frac{1}{2} \ln(1+x^2) + C.
\]
Therefore,
\[
\int \arctan x \, dx = x \arctan x - \frac{1}{2} \ln(1+x^2) + C.
\]
\end{example}
\begin{example}
Evaluate $\displaystyle \int x \ln x \, dx$.
Let $u = \ln x$ and $dv = x \, dx$. Then $du = \frac{1}{x} \, dx$ and $v = \frac{x^2}{2}$. Thus,
\[
\int x \ln x \, dx = \frac{x^2}{2} \ln x - \int \frac{x^2}{2} \cdot \frac{1}{x} \, dx = \frac{x^2}{2} \ln x - \frac{1}{2} \int x \, dx = \frac{x^2}{2} \ln x - \frac{x^2}{4} + C.
\]
\end{example}
\section{Special Cases of Integration}
\begin{theorem}[Reduction Formula for $\int \sin^n x \, dx$]
For $n \geq 2$,
\[
\int \sin^n x \, dx = -\frac{1}{n} \sin^{n-1} x \cos x + \frac{n-1}{n} \int \sin^{n-2} x \, dx.
\]
\end{theorem}
\begin{proof}
Write $\sin^n x = \sin^{n-1} x \sin x$ and integrate by parts with $u = \sin^{n-1} x$ and $dv = \sin x \, dx$. Then $du = (n-1)\sin^{n-2} x \cos x \, dx$ and $v = -\cos x$. Thus,
\[
\int \sin^n x \, dx = -\sin^{n-1} x \cos x + (n-1) \int \sin^{n-2} x \cos^2 x \, dx.
\]
Using $\cos^2 x = 1 - \sin^2 x$,
\[
\int \sin^n x \, dx = -\sin^{n-1} x \cos x + (n-1) \int \sin^{n-2} x \, dx - (n-1) \int \sin^n x \, dx.
\]
Adding $(n-1)\int \sin^n x \, dx$ to both sides gives
\[
n \int \sin^n x \, dx = -\sin^{n-1} x \cos x + (n-1) \int \sin^{n-2} x \, dx.
\]
Dividing by $n$ yields the formula.
\end{proof}
\begin{theorem}[Reduction Formula for $\int \cos^n x \, dx$]
For $n \geq 2$,
\[
\int \cos^n x \, dx = \frac{1}{n} \cos^{n-1} x \sin x + \frac{n-1}{n} \int \cos^{n-2} x \, dx.
\]
\end{theorem}
\begin{proof}
Similar to the sine reduction formula, using $u = \cos^{n-1} x$ and $dv = \cos x \, dx$.
\end{proof}
\begin{example}
Evaluate $\displaystyle \int \sin^4 x \, dx$.
Using the reduction formula with $n = 4$,
\[
\int \sin^4 x \, dx = -\frac{1}{4} \sin^3 x \cos x + \frac{3}{4} \int \sin^2 x \, dx.
\]
For $\int \sin^2 x \, dx$, use the identity $\sin^2 x = \frac{1 - \cos 2x}{2}$:
\[
\int \sin^2 x \, dx = \frac{x}{2} - \frac{\sin 2x}{4} + C.
\]
Therefore,
\[
\int \sin^4 x \, dx = -\frac{1}{4} \sin^3 x \cos x + \frac{3}{4} \left( \frac{x}{2} - \frac{\sin 2x}{4} \right) + C = -\frac{1}{4} \sin^3 x \cos x + \frac{3x}{8} - \frac{3 \sin 2x}{16} + C.
\]
\end{example}
\begin{theorem}[Integrals of the Form $\int \tan^n x \, dx$ and $\int \sec^n x \, dx$]
For $n \geq 2$,
\[
\int \tan^n x \, dx = \frac{\tan^{n-1} x}{n-1} - \int \tan^{n-2} x \, dx.
\]
For $n \geq 2$,
\[
\int \sec^n x \, dx = \frac{\sec^{n-2} x \tan x}{n-1} + \frac{n-2}{n-1} \int \sec^{n-2} x \, dx.
\]
\end{theorem}
\begin{proof}
For the tangent formula, write $\tan^n x = \tan^{n-2} x \tan^2 x = \tan^{n-2} x (\sec^2 x - 1)$. Then
\[
\int \tan^n x \, dx = \int \tan^{n-2} x \sec^2 x \, dx - \int \tan^{n-2} x \, dx.
\]
The first integral is $\frac{\tan^{n-1} x}{n-1}$ by substitution $u = \tan x$. For the secant formula, integrate by parts with $u = \sec^{n-2} x$ and $dv = \sec^2 x \, dx$.
\end{proof}
\begin{example}
Evaluate $\displaystyle \int \tan^3 x \, dx$.
Using the formula with $n = 3$,
\[
\int \tan^3 x \, dx = \frac{\tan^2 x}{2} - \int \tan x \, dx = \frac{\tan^2 x}{2} + \ln|\cos x| + C.
\]
\end{example}
\begin{theorem}[Integrals of the Form $\int \sin^m x \cos^n x \, dx$]
\begin{enumerate}
\item If $m$ is odd, substitute $u = \cos x$.
\item If $n$ is odd, substitute $u = \sin x$.
\item If both $m$ and $n$ are even, use the half-angle identities
\[
\sin^2 x = \frac{1 - \cos 2x}{2}, \quad \cos^2 x = \frac{1 + \cos 2x}{2}.
\]
\end{enumerate}
\end{theorem}
\begin{example}
Evaluate $\displaystyle \int \sin^3 x \cos^2 x \, dx$.
Since $m = 3$ is odd, let $u = \cos x$. Then $du = -\sin x \, dx$ and $\sin^2 x = 1 - \cos^2 x = 1 - u^2$. Thus,
\[
\int \sin^3 x \cos^2 x \, dx = \int \sin^2 x \cos^2 x \sin x \, dx = \int (1 - u^2) u^2 (-du) = \int (u^4 - u^2) \, du = \frac{u^5}{5} - \frac{u^3}{3} + C = \frac{\cos^5 x}{5} - \frac{\cos^3 x}{3} + C.
\]
\end{example}
\begin{example}
Evaluate $\displaystyle \int \sin^2 x \cos^4 x \, dx$.
Both exponents are even. Use $\sin^2 x = \frac{1 - \cos 2x}{2}$ and $\cos^2 x = \frac{1 + \cos 2x}{2}$:
\[
\sin^2 x \cos^4 x = \left( \frac{1 - \cos 2x}{2} \right) \left( \frac{1 + \cos 2x}{2} \right)^2 = \frac{1}{8} (1 - \cos 2x)(1 + 2\cos 2x + \cos^2 2x).
\]
Expanding,
\[
= \frac{1}{8} (1 + 2\cos 2x + \cos^2 2x - \cos 2x - 2\cos^2 2x - \cos^3 2x) = \frac{1}{8} (1 + \cos 2x - \cos^2 2x - \cos^3 2x).
\]
Using $\cos^2 2x = \frac{1 + \cos 4x}{2}$,
\[
= \frac{1}{8} \left( 1 + \cos 2x - \frac{1 + \cos 4x}{2} - \cos^3 2x \right) = \frac{1}{16} (1 + 2\cos 2x - \cos 4x - 2\cos^3 2x).
\]
Thus,
\[
\int \sin^2 x \cos^4 x \, dx = \frac{1}{16} \left( x + \sin 2x - \frac{\sin 4x}{4} - 2 \int \cos^3 2x \, dx \right).
\]
For $\int \cos^3 2x \, dx$, let $u = \sin 2x$, $du = 2\cos 2x \, dx$, and $\cos^2 2x = 1 - u^2$:
\[
\int \cos^3 2x \, dx = \int \cos^2 2x \cos 2x \, dx = \frac{1}{2} \int (1 - u^2) \, du = \frac{1}{2} \left( u - \frac{u^3}{3} \right) = \frac{\sin 2x}{2} - \frac{\sin^3 2x}{6}.
\]
Substituting back yields the final answer.
\end{example}
\begin{theorem}[Integrals of the Form $\int \sin(mx) \cos(nx) \, dx$]
Use the product-to-sum identities:
\[
\sin A \cos B = \frac{1}{2} [\sin(A-B) + \sin(A+B)],
\]
\[
\sin A \sin B = \frac{1}{2} [\cos(A-B) - \cos(A+B)],
\]
\[
\cos A \cos B = \frac{1}{2} [\cos(A-B) + \cos(A+B)].
\]
\end{theorem}
\begin{example}
Evaluate $\displaystyle \int \sin 3x \cos 5x \, dx$.
Using the identity,
\[
\sin 3x \cos 5x = \frac{1}{2} [\sin(3x-5x) + \sin(3x+5x)] = \frac{1}{2} [\sin(-2x) + \sin 8x] = \frac{1}{2} [-\sin 2x + \sin 8x].
\]
Thus,
\[
\int \sin 3x \cos 5x \, dx = \frac{1}{2} \left( \frac{\cos 2x}{2} - \frac{\cos 8x}{8} \right) + C = \frac{\cos 2x}{4} - \frac{\cos 8x}{16} + C.
\]
\end{example}
\section{Integrals of Rational Functions by Partial Fractions}
\begin{definition}[Rational Function]
A \textit{rational function} is a function of the form $f(x) = \frac{P(x)}{Q(x)}$, where $P$ and $Q$ are polynomials and $Q(x) \neq 0$.
\end{definition}
\begin{definition}[Proper and Improper Rational Functions]
A rational function $\frac{P(x)}{Q(x)}$ is \textit{proper} if the degree of $P$ is less than the degree of $Q$. It is \textit{improper} if the degree of $P$ is greater than or equal to the degree of $Q$.
\end{definition}
\begin{theorem}[Partial Fraction Decomposition]
Any proper rational function $\frac{P(x)}{Q(x)}$ can be expressed as a sum of simpler fractions, provided $Q(x)$ is factored completely into linear and irreducible quadratic factors. The form of the decomposition depends on the factors of $Q(x)$:
\begin{enumerate}
\item For each distinct linear factor $(ax + b)$, include a term $\frac{A}{ax + b}$.
\item For each repeated linear factor $(ax + b)^k$, include terms $\frac{A_1}{ax + b} + \frac{A_2}{(ax + b)^2} + \dots + \frac{A_k}{(ax + b)^k}$.
\item For each distinct irreducible quadratic factor $(ax^2 + bx + c)$, include a term $\frac{Ax + B}{ax^2 + bx + c}$.
\item For each repeated irreducible quadratic factor $(ax^2 + bx + c)^k$, include terms $\frac{A_1 x + B_1}{ax^2 + bx + c} + \dots + \frac{A_k x + B_k}{(ax^2 + bx + c)^k}$.
\end{enumerate}
\end{theorem}
\begin{example}
Evaluate $\displaystyle \int \frac{1}{x^2 - 1} \, dx$.
Factor the denominator: $x^2 - 1 = (x - 1)(x + 1)$. Write
\[
\frac{1}{(x-1)(x+1)} = \frac{A}{x-1} + \frac{B}{x+1}.
\]
Multiply by $(x-1)(x+1)$: $1 = A(x+1) + B(x-1)$.
Set $x = 1$: $1 = 2A \implies A = \frac{1}{2}$.
Set $x = -1$: $1 = -2B \implies B = -\frac{1}{2}$.
Thus,
\[
\int \frac{1}{x^2 - 1} \, dx = \frac{1}{2} \int \frac{1}{x-1} \, dx - \frac{1}{2} \int \frac{1}{x+1} \, dx = \frac{1}{2} \ln|x-1| - \frac{1}{2} \ln|x+1| + C = \frac{1}{2} \ln\left| \frac{x-1}{x+1} \right| + C.
\]
\end{example}
\begin{example}
Evaluate $\displaystyle \int \frac{2x + 3}{x^2 + 3x + 2} \, dx$.
Factor: $x^2 + 3x + 2 = (x+1)(x+2)$. Write
\[
\frac{2x + 3}{(x+1)(x+2)} = \frac{A}{x+1} + \frac{B}{x+2}.
\]
Multiply: $2x + 3 = A(x+2) + B(x+1)$.
Set $x = -1$: $1 = A(1) \implies A = 1$.
Set $x = -2$: $-1 = B(-1) \implies B = 1$.
Thus,
\[
\int \frac{2x + 3}{x^2 + 3x + 2} \, dx = \int \frac{1}{x+1} \, dx + \int \frac{1}{x+2} \, dx = \ln|x+1| + \ln|x+2| + C = \ln|(x+1)(x+2)| + C.
\]
\end{example}
\begin{example}
Evaluate $\displaystyle \int \frac{x^2 + 2x - 1}{x^3 - x} \, dx$.
Factor: $x^3 - x = x(x-1)(x+1)$. Write
\[
\frac{x^2 + 2x - 1}{x(x-1)(x+1)} = \frac{A}{x} + \frac{B}{x-1} + \frac{C}{x+1}.
\]
Multiply: $x^2 + 2x - 1 = A(x-1)(x+1) + Bx(x+1) + Cx(x-1)$.
Set $x = 0$: $-1 = A(-1)(1) \implies A = 1$.
Set $x = 1$: $1 + 2 - 1 = 2 = B(1)(2) \implies B = 1$.
Set $x = -1$: $1 - 2 - 1 = -2 = C(-1)(-2) \implies C = -1$.
Thus,
\[
\int \frac{x^2 + 2x - 1}{x^3 - x} \, dx = \int \frac{1}{x} \, dx + \int \frac{1}{x-1} \, dx - \int \frac{1}{x+1} \, dx = \ln|x| + \ln|x-1| - \ln|x+1| + C = \ln\left| \frac{x(x-1)}{x+1} \right| + C.
\]
\end{example}
\begin{example}
Evaluate $\displaystyle \int \frac{x^2 + 1}{x(x^2 + 1)} \, dx$ (note: this simplifies before decomposition).
Actually, $\frac{x^2 + 1}{x(x^2 + 1)} = \frac{1}{x}$ for $x \neq 0$. Thus,
\[
\int \frac{x^2 + 1}{x(x^2 + 1)} \, dx = \int \frac{1}{x} \, dx = \ln|x| + C.
\]
\end{example}
\begin{example}
Evaluate $\displaystyle \int \frac{1}{x^2(x+1)} \, dx$.
The factor $x^2$ is a repeated linear factor. Write
\[
\frac{1}{x^2(x+1)} = \frac{A}{x} + \frac{B}{x^2} + \frac{C}{x+1}.
\]
Multiply: $1 = Ax(x+1) + B(x+1) + Cx^2$.
Set $x = 0$: $1 = B(1) \implies B = 1$.
Set $x = -1$: $1 = C(1) \implies C = 1$.
Expand and compare coefficients: $1 = A(x^2 + x) + (x+1) + x^2 = (A+1)x^2 + (A+1)x + 1$.
Thus $A+1 = 0 \implies A = -1$.
Therefore,
\[
\int \frac{1}{x^2(x+1)} \, dx = -\int \frac{1}{x} \, dx + \int \frac{1}{x^2} \, dx + \int \frac{1}{x+1} \, dx = -\ln|x| - \frac{1}{x} + \ln|x+1| + C = \ln\left| \frac{x+1}{x} \right| - \frac{1}{x} + C.
\]
\end{example}
\begin{example}
Evaluate $\displaystyle \int \frac{2x + 1}{x^2 + 2x + 5} \, dx$.
The denominator $x^2 + 2x + 5$ is irreducible (discriminant $4 - 20 = -16 < 0$). Complete the square: $x^2 + 2x + 5 = (x+1)^2 + 4$. Write the numerator as $2x + 1 = 2(x+1) - 1$. Thus,
\[
\int \frac{2x + 1}{x^2 + 2x + 5} \, dx = \int \frac{2(x+1) - 1}{(x+1)^2 + 4} \, dx = \int \frac{2(x+1)}{(x+1)^2 + 4} \, dx - \int \frac{1}{(x+1)^2 + 4} \, dx.
\]
For the first integral, let $u = (x+1)^2 + 4$, $du = 2(x+1) \, dx$:
\[
\int \frac{2(x+1)}{(x+1)^2 + 4} \, dx = \ln((x+1)^2 + 4) + C_1.
\]
For the second, let $v = \frac{x+1}{2}$, $dv = \frac{1}{2} dx$:
\[
\int \frac{1}{(x+1)^2 + 4} \, dx = \frac{1}{2} \arctan\left( \frac{x+1}{2} \right) + C_2.
\]
Therefore,
\[
\int \frac{2x + 1}{x^2 + 2x + 5} \, dx = \ln(x^2 + 2x + 5) - \frac{1}{2} \arctan\left( \frac{x+1}{2} \right) + C.
\]
\end{example}
\begin{example}
Evaluate $\displaystyle \int \frac{x^3 + x}{x - 1} \, dx$.
The rational function is improper. Perform polynomial long division:
\[
\frac{x^3 + x}{x - 1} = x^2 + x + 2 + \frac{2}{x - 1}.
\]
Thus,
\[
\int \frac{x^3 + x}{x - 1} \, dx = \int \left( x^2 + x + 2 + \frac{2}{x-1} \right) \, dx = \frac{x^3}{3} + \frac{x^2}{2} + 2x + 2\ln|x-1| + C.
\]
\end{example}
\begin{theorem}[General Method for Integrating Rational Functions]
\begin{enumerate}
\item If the rational function is improper, perform polynomial long division.
\item Factor the denominator completely into linear and irreducible quadratic factors.
\item Decompose the proper rational function into partial fractions.
\item Integrate each partial fraction term using standard formulas.
\end{enumerate}
\end{theorem}
\begin{example}
Evaluate $\displaystyle \int \frac{4x^2 + 3x + 2}{x^3 + x^2} \, dx$.
Factor: $x^3 + x^2 = x^2(x+1)$. Write
\[
\frac{4x^2 + 3x + 2}{x^2(x+1)} = \frac{A}{x} + \frac{B}{x^2} + \frac{C}{x+1}.
\]
Multiply: $4x^2 + 3x + 2 = Ax(x+1) + B(x+1) + Cx^2$.
Set $x = 0$: $2 = B \implies B = 2$.
Set $x = -1$: $4 - 3 + 2 = 3 = C(1) \implies C = 3$.
Expand: $4x^2 + 3x + 2 = A(x^2 + x) + 2x + 2 + 3x^2 = (A+3)x^2 + (A+2)x + 2$.
Comparing coefficients: $A + 3 = 4 \implies A = 1$.
Thus,
\[
\int \frac{4x^2 + 3x + 2}{x^3 + x^2} \, dx = \int \frac{1}{x} \, dx + 2 \int \frac{1}{x^2} \, dx + 3 \int \frac{1}{x+1} \, dx = \ln|x| - \frac{2}{x} + 3\ln|x+1| + C.
\]
\end{example}
\begin{example}
Evaluate $\displaystyle \int \frac{x^2 - 2x - 3}{(x-1)(x^2+1)} \, dx$.
Write
\[
\frac{x^2 - 2x - 3}{(x-1)(x^2+1)} = \frac{A}{x-1} + \frac{Bx + C}{x^2 + 1}.
\]
Multiply: $x^2 - 2x - 3 = A(x^2 + 1) + (Bx + C)(x - 1)$.
Set $x = 1$: $1 - 2 - 3 = -4 = A(2) \implies A = -2$.
Expand: $x^2 - 2x - 3 = -2x^2 - 2 + Bx^2 - Bx + Cx - C = (B - 2)x^2 + (C - B)x + (-2 - C)$.
Comparing coefficients:
$B - 2 = 1 \implies B = 3$.
$C - B = -2 \implies C - 3 = -2 \implies C = 1$.
Check constant: $-2 - C = -3 \implies -2 - 1 = -3$, correct.
Thus,
\[
\int \frac{x^2 - 2x - 3}{(x-1)(x^2+1)} \, dx = -2 \int \frac{1}{x-1} \, dx + \int \frac{3x + 1}{x^2 + 1} \, dx.
\]
For the second integral,
\[
\int \frac{3x + 1}{x^2 + 1} \, dx = \frac{3}{2} \int \frac{2x}{x^2 + 1} \, dx + \int \frac{1}{x^2 + 1} \, dx = \frac{3}{2} \ln(x^2 + 1) + \arctan x + C.
\]
Therefore,
\[
\int \frac{x^2 - 2x - 3}{(x-1)(x^2+1)} \, dx = -2\ln|x-1| + \frac{3}{2} \ln(x^2 + 1) + \arctan x + C.
\]
\end{example}
\end{document}
